\documentclass[10pt,a4paper]{amsart}
\usepackage[margin=19mm]{geometry}
\usepackage{amsmath,amssymb,mathtools}
\usepackage[scr=rsfs,cal=euler]{mathalfa}
\usepackage{xfrac}
\usepackage[unicode,hidelinks,bookmarks=false]{hyperref}
\newcommand{\ie}{i.e.\ }
\newcommand{\cf}{cf.\ }
\newcommand{\resp}{respectively\ }
\newcommand{\Subsection}{\subsection}
\providecommand{\nobreakdash}{\nobreak\hbox{-}}
\setlength{\emergencystretch}{2em}
\multlinegap0pt
\usepackage{bm}
\usepackage{bbm}
\usepackage{dsfont}
\usepackage{xspace}
\usepackage{cancel}
\usepackage{mathtools}
\usepackage{amsthm}
\makeatletter
\@ifundefined{theorem}{\newtheorem{theorem}{Theorem}}{}
\@ifundefined{corollary}{\newtheorem{corollary}[theorem]{Corollary}}{}
\makeatother
\title[A small noise approximation for Muller's Ratchet]{A small noise approximation for Muller's Ratchet}
\author{Carola Sophia Heinzel, Peter Pfaffelhuber, Anton Wakolbinger}
\date{Source: arXiv:2606.15842v1 (CC BY 4.0)}
\begin{document}
\maketitle

\noindent\emph{Source and licence.} A small noise approximation for Muller's Ratchet. Version arXiv:2606.15842v1 (CC BY 4.0), distributed under the Creative Commons Attribution 4.0 International licence, as recorded in the arXiv licence metadata for this exact version and shown on the abstract page https://arxiv.org/abs/2606.15842v1. Full rights notice: licenses/arxiv-2606.15842-CC-BY-4.0.txt.\par\smallskip

\noindent\textit{Editorial scope.} Counted unit: the Corollary whose source label is \texttt{cor2.4} in the article (source0.tex lines 590--617, source label \texttt{cor2.4}), delivered together with the article's own complete original proof of it (source0.tex lines 658--760, one \texttt{proof} environment of 103 lines). No mathematics is written, repaired or re-derived: statement and proof are the article's own text line for line. The proof is detached from the statement in the source: the article states the result at lines 590--617 and prints its proof at lines 658--760, and the proof environment's own header names the statement, so the association is the article's own. Required context retained uncounted: eq:dynsys1 (lines 509--512). Every definition, display and label of those lines is retained unchanged. Omitted from this package and not redistributed: 1--508, 513--589, 618--657, 761--1370. Nothing inside a retained excerpt is omitted or altered: every retained body is delivered exactly as the article's lines are written, so no omission receipt of any kind accompanies this package. Citations: the retained excerpts carry no citation command at all, so no bibliography entry of the article is transcribed and no reference is redistributed. Reference adaptation: none was needed and none was made; every reference inside the delivered unit resolves inside it, so no unresolved reference or citation is produced. Printed slips kept literal: no printed word, symbol, hypothesis or number of the counted statement or of its proof was altered. Layout and class: the article's own class is \texttt{imsart} with packages including amsthm, amsmath, amsfonts, amssymb, natbib, hyperref, graphicx, float, xcolor, rotating, caption, enumitem, ulem, verbatim; the wrapper is the lane's pinned \texttt{amsart} editorial wrapper at 10pt on A4, which restates the theorem-like environments the retained text needs and adds the article's own macro definitions that the retained text needs (none), with extra wrapper packages (bm, bbm, dsfont, xspace, cancel, mathtools). The delivered numbering is the unit's own. Assets: no retained span contains a figure environment, a graphics-inclusion command, a PDF-page inclusion, a table or a TikZ picture; no image file is copied into this package, the package declares no asset dependency and no image file is redistributed with it. Licence: version arXiv:2606.15842v1 of this article is distributed under the Creative Commons Attribution 4.0 International licence, as recorded in the arXiv licence metadata for this exact version and shown on the abstract page https://arxiv.org/abs/2606.15842v1. A full rights notice is delivered at licenses/arxiv-2606.15842-CC-BY-4.0.txt. Characters: every retained excerpt is pure ASCII, so the delivered engine, XeLaTeX with its default Latin Modern text font, reports no missing character.
\begin{align}
  \label{eq:dynsys1}
  \dot x_k & = \lambda(x_{k-1}-x_k) + \alpha \sum_{j=0}^\infty (j - k) x_j x_k,
\end{align}

\begin{corollary}\label{cor2.4}
  Let $(S_t)_{t\geq 0}$ be the semigroup generated by $G^{(1)}$. Started in $x \in \mathbb S$, the solution of \eqref{eq:dynsys1} satisfies, for $\xi \in \mathbb R_+$, 
\begin{equation}
\label{eq:sol2}
\begin{aligned}
  S_t g_{\xi}(x) & = g_{\xi}(x(t)) = \frac{g_{\alpha t + \xi}(x)}{g_{\alpha t}(x)}
  \exp\Big(-\theta(1-e^{- \xi})(1-e^{-\alpha t})\Big)
\end{aligned}
\end{equation}
as well as, for $\xi, \eta \in \mathbb R_+$,
\begin{align}
    \label{eq:sol3}
    G^{(1)}S_t g_{\xi} & = \exp\Big(-\theta(1-e^{- \xi})(1-e^{-\alpha t})\Big)\Big(\lambda e^{-\alpha t}(1 - e^{-\xi})\frac{g_{\alpha  t + \xi}}{g_{\alpha t}} + \frac{\partial}{\partial t} \frac{g_{\alpha  t + \xi}}{g_{\alpha t}}\Big), \\ \label{eq:sol32}
    G^{(1)}S_t g_{\xi} g_{\eta} & = (S_t g_{\xi})\, G^{(1)}S_t g_{\eta} + (S_t g_{ \eta})\, G^{(1)}S_t g_{ \xi}, \\ 
    \label{eq:sol4}
    \nu \cdot G^{(2)}S_t g_{ \xi} & = 
    \frac{g_{\alpha  t + \xi}(x)g_{2\alpha t} - g_{\alpha t} g_{2\alpha t+\xi}}{g_{\alpha t}^3}\exp\Big(-\theta(1-e^{-\xi})(1-e^{-\alpha t})\Big), \\
    \nu \cdot G^{(2)}S_t g_{\xi} g_{\eta}  & = \exp\Big(-\theta(2-e^{-\xi} - e^{- \eta})(1-e^{-\alpha t})\Big) \notag \\ 
    \label{eq:sol42} &  \cdot 
    \frac{
    3g_{\alpha t + \xi} g_{\alpha  t + \eta}g_{2\alpha t} - 2 g_{\alpha  t + \xi} g_{\alpha t}g_{2\alpha t + \eta} - 2 g_{\alpha t + \eta} g_{\alpha t} g_{2\alpha t+\xi} + g_{\alpha t}^2 g_{2\alpha t+\xi+\eta}}{g_{\alpha t}^4}. 
\end{align}
Moreover, for $f_{t,\xi,\eta} := \nu \cdot G^{(2)}S_t g_{\xi} g_{\eta}$, there are bounded functions $g_{s, t, \xi, \eta}, g^1_{s, t, \xi, \eta}, g^2_{s, t, \xi, \eta} \in \text{span} \{g_{\xi_1} \cdots g_{\xi_m}: m \in \mathbb N, \xi_1,...,\xi_m \in \mathbb R_+ \}$ and $a, a_1, a_2$ (all not depending on $\nu$) such that 
\begin{align}
    \label{eq:morebounds}
    S_s f_{t,\xi, \eta} = \frac{g_{s,t,\xi, \eta}}{g_{\alpha (s+t)}^{a}}, \qquad G^{(1)} S_s f_{t,\xi, \eta} = \frac{g^1_{s,t,\xi, \eta}}{g_{\alpha (s+t)}^{a_1}}, \qquad \nu \cdot G^{(2)} S_s f_{t,\xi, \eta} = \frac{g^2_{s,t,\xi, \eta}}{g_{\alpha (s+t)}^{a_2}}.
\end{align}
\end{corollary}

\begin{proof}[Proof of Corollary~\ref{cor2.4}]
  For \eqref{eq:sol2}, we compute
  \begin{align*}
    g_{ \xi}(x(t)) & = \sum_{k=0}^\infty e^{- \xi k} x_k(t) \\ & = \frac{1}{g_{\alpha t}(x)} \! \sum_{i=0}^\infty x_i e^{-i(\alpha t+\xi)} \Big(\underbrace{ \sum_{k=i}^\infty  \frac{\theta^{k-i} e^{- (k-i)\xi}}{(k-i)!}
      (1-e^{-\alpha t})^{k-i}}_{= \exp\big( \theta(e^{- \xi} ( 1\! - e^{-\alpha t}) \big) }\Big)\Big)\exp\Big(\!-\theta(1 \! -e^{-\alpha t})\Big)\\ & =
    \frac{g_{\alpha t + \xi}(x)}{ g_{\alpha t}(x)}
    \exp\Big(-\theta(1-e^{- \xi})(1-e^{-\alpha t})\Big).
  \end{align*}
  Then,
  \begin{align*}
    G^{(1)}S_t g_{\xi}(x) & = \frac{\partial}{\partial t}S_t g_{\xi}(x) \\ & = \exp\Big(-\theta(1-e^{- \xi})(1-e^{-\alpha t})\Big)\Big(\lambda e^{-\alpha t}(1 - e^{- \xi})\frac{g_{\alpha  t + \xi}(x)}{g_{\alpha t}(x)} + \frac{\partial}{\partial t} \frac{g_{\alpha t + \xi}(x)}{g_{\alpha t}(x)}\Big),
\end{align*}
which shows \eqref{eq:sol3}. From this, \eqref{eq:sol32} follows since $S_t(g_{\xi} g_{\eta}) = (S_t g_{\xi})(S_tg_{\eta})$ and $G^{(1)}$ is first order. 

Next, for $G^{(2)}S_t g_{\xi}$, taking a look at the form of $S_tg_{\xi}(x)$, we need to compute
\begin{align*}
  \sum_{k,\ell=0}^\infty & 
  x_k(\delta_{k\ell} - x_\ell) \frac{\partial^2}{\partial x_k
    \partial x_\ell} \frac{g_{\alpha t + \xi}(x)}{g_{\alpha t}(x)} 
  \\
  & = \sum_{k,\ell=0}^\infty
  x_k(\delta_{k\ell} - x_\ell) \frac{\partial}{\partial x_k} \frac{g_{\alpha t}(x) e^{-(\alpha t + \xi)\ell} - g_{\alpha t + \xi}(x) e^{-\alpha t \ell}}{g_{\alpha t}(x)^2} 
  \\ & = \sum_{k,\ell=0}^\infty
  x_k(\delta_{k\ell} - x_\ell) \cdot \big(g_{\alpha t}(x) ( e^{-\alpha t k} e^{-(\alpha t + \xi)\ell} - e^{-(\alpha t + \xi) k} e^{-\alpha t \ell}) \\ & \qquad \qquad \qquad \qquad \qquad - 2 e^{-\alpha tk}(g_{\alpha t}(x) e^{-(\alpha t + \xi)\ell} - g_{\alpha t + \xi}(x) e^{-\alpha t \ell})\big) / g_{\alpha t}(x)^3
  \\ & =  
  -2\frac{g_{\alpha t}(x)g_{2\alpha t+\xi}(x) - g_{\alpha t +\xi}(x)g_{2\alpha t}(x)}{g_{\alpha t}(x)^3}.
\end{align*}
Therefore, 
\begin{align*}
    \nu G^{(2)} S_{t} g_{ \xi}(x) & = \frac{1}{2} \sum_{k,\ell = 0}^\infty x_k ( \delta_{kl} - x_\ell) \frac{\partial^2 S_{t} g_{ \xi} (x)}{\partial x_k \partial x_\ell} \\ & = 
    \frac{1}{2} \sum_{k,\ell = 0}^\infty x_k ( \delta_{kl} - x_\ell) \frac{\partial^2 }{\partial x_k \partial x_\ell} \frac{g_{\alpha t +\xi}(x)}{g_{\alpha t}(x)}
    \exp\Big(-\theta(1-e^{-\xi})(1-e^{-\alpha t})\Big) \\ & = 
    \frac{g_{\alpha t + \xi}(x)g_{\alpha 2t}(x) - g_{\alpha t}(x) g_{2\alpha t+\xi}(x)}{g_{\alpha t}(x)^3}\exp\Big(-\theta(1-e^{- \xi})(1-e^{-\alpha t})\Big).
\end{align*}
Let us turn to \eqref{eq:sol42}. We calculate, omitting the $x$-dependence,
\begin{align*}
  \sum_{k,\ell=0}^\infty & 
  x_k(\delta_{k\ell} - x_\ell) \Big(\frac{\partial}{\partial x_k}\frac{g_{\alpha  t + \xi}}{g_{\alpha t}} \Big) \Big(\frac{\partial}{\partial x_\ell} \frac{g_{\alpha t + \eta}}{g_{\alpha t}}\Big)
  \\
  & = \sum_{k,\ell=0}^\infty
  x_k(\delta_{k\ell} - x_\ell) \frac{(g_{\alpha t} e^{-(\alpha t + \xi)k} - g_{\alpha t + \xi} e^{-\alpha t k})(g_{\alpha t} e^{-(\alpha t + \eta)\ell} - g_{\alpha t + \eta} e^{-\alpha t \ell})}{g_{\alpha t}^4} 
  \\ & = \Big(g_{\alpha t}^2 g_{2\alpha t+\xi+\eta} - g_{\alpha t} g_{2\alpha t+\xi} g_{\alpha t + \eta} - g_{\alpha t} g_{2\alpha t+\eta}g_{\alpha t + \xi} + g_{\alpha t +\xi}g_{\alpha t + \eta}g_{2\alpha t})\Big) / g_{\alpha t}^4.
\end{align*}
Therefore, since $S_t g_{ \xi} g_{ \eta} = (S_t g_{ \xi})(S_t g_{ \eta})$, \eqref{eq:sol42} follows from
\begin{align*}
    \exp\Big( & \theta(2 - e^{- \xi} - e^{- \eta})(1 - e^{-\alpha t}) \Big) \nu G^{(2)}S_t g_{ \xi} g_{ \eta}  \\ & = \nu\frac{g_{\alpha t + \xi}}{g_{\alpha t}} G^{(2)} \frac{g_{\alpha t + \eta}}{g_{\alpha t}} + \nu\frac{g_{\alpha t + \eta}}{g_{\alpha t}} G^{(2)} \frac{g_{\alpha t + \xi}}{g_{\alpha t}} \\ & \qquad \qquad \qquad \qquad + \sum_{k,\ell=0}^\infty 
    x_k(\delta_{k\ell} - x_\ell) \Big(\frac{\partial}{\partial x_k}\frac{g_{\alpha t + \xi}}{g_{\alpha t}} \Big) \Big(\frac{\partial}{\partial x_\ell} \frac{g_{\alpha t + \eta}}{g_{\alpha t}}\Big)
    \\ & = 
    \Big(g_{\alpha t + \xi}(g_{\alpha t + \eta}g_{\alpha 2t} - g_{\alpha t}g_{2\alpha t + \eta}) + g_{\alpha t + \eta} (g_{\alpha t + \xi}g_{2\alpha t} - g_{\alpha t}g_{2\alpha t+\xi}) \\ & \qquad + g_{\alpha t}^2 g_{2\alpha t+\xi+\eta} - g_{\alpha t} g_{2\alpha t+\xi}g_{\alpha t + \eta} - g_{\alpha t} g_{2\alpha t+\eta}g_{\alpha t + \xi} + g_{\alpha t +\xi}g_{\alpha t + \eta}g_{2\alpha t}\Big)/g_{\alpha t}^4 \\ & = \Big(
    2g_{\alpha t + \xi}(g_{\alpha t + \eta}g_{2\alpha t} - g_{\alpha t}g_{2\alpha t + \eta}) + 2g_{\alpha t + \eta}(g_{\alpha t + \xi}g_{2\alpha t} - g_{\alpha t}g_{2\alpha t+\xi}) \\ & \qquad \qquad \qquad \qquad \qquad \qquad \qquad \qquad + g_{\alpha t}^2 g_{2\alpha t+\xi+\eta} - g_{\alpha t + \xi}g_{\alpha t + \eta}g_{2\alpha t}\Big) / g_{\alpha t}^4.
\end{align*}
Let us prove \eqref{eq:morebounds}. Define
$$
c(s,\xi):=
\exp\Big(
-\theta(1-e^{-\xi})(1-e^{-\alpha s})
\Big).
$$
Then, it holds
$$
S_s g_{\xi}  =c(s,\xi)\frac{g_{\alpha s+\xi}  }{g_{\alpha s}  },
$$
which leads to
\begin{align*}
S_s\left(
\frac{
g_{\xi}g_{\eta}g_{\zeta}
}{g_{\alpha t}^4}
\right)  
&=
\frac{
S_s g_{\xi}  S_s g_{\eta}  S_s g_{\zeta}  
}{
(S_s g_{\alpha t}  )^4
} \\
&=
\frac{
c(s,\xi)c(s,\eta)c(s,\zeta)
}{
c(s,\alpha t)^4
}
\frac{
\frac{g_{\alpha s + \xi}  }{g_{\alpha s}  }
\frac{g_{\alpha s + \eta}  }{g_{\alpha s}  }
\frac{g_{\alpha s + \zeta}  }{g_{\alpha s}  }
}{
\left(\frac{g_{\alpha (s+t)}  }{g_{\alpha s}  }\right)^4
} \\
&=
\frac{
c(s,\xi)c(s,\eta)c(s,\zeta)
}{
c(s,\alpha t)^4
}
\frac{
g_{\alpha s}  g_{\alpha s + \xi}  g_{\alpha s + \eta}  g_{\alpha s + \zeta}  
}{
g_{\alpha (s+t)}  ^4
}.
\end{align*}

The rest of the proof of \eqref{eq:morebounds} follows along the same lines.
\end{proof}

\end{document}
